% !TeX program = xelatex
% Working manuscript: Q1 integrated from teammate upload; Q2 is a draft.
\documentclass[UTF8,a4paper,zihao=-4,fontset=fandol]{ctexart}
\usepackage[margin=25mm,headheight=15pt]{geometry}
\usepackage{amsmath,amssymb,mathtools}
\usepackage{graphicx,booktabs,tabularx,array}
\usepackage{enumitem,needspace,tcolorbox,fancyhdr}
\usepackage{tikz}
\usepackage{caption}
\captionsetup{font=small,labelfont=bf,labelsep=quad,hypcap=false}
\usepackage{tikz}
\usetikzlibrary{arrows.meta,calc,patterns}
\definecolor{qoneink}{HTML}{203045}
\definecolor{qoneblue}{HTML}{347DB0}
\definecolor{qoneorange}{HTML}{CB8435}
\definecolor{qonepurple}{HTML}{8A61A8}
\definecolor{qoneregion}{HTML}{207C83}
\definecolor{qonefail}{HTML}{C94E48}
\definecolor{qonecover}{HTML}{348664}

\newcommand{\QOneCaseFigure}{\resizebox{\linewidth}{!}{%
\begin{tikzpicture}[>=Stealth,line cap=round,line join=round,
  font=\small,text=qoneink,
  dot/.style={circle,inner sep=1.6pt,fill=qoneink},
  info/.style={fill=white,inner sep=2pt,align=center}]
\path[use as bounding box] (0,.64) rectangle (23,9.2);
\fill[white] (0,.64) rectangle (23,9.2);
\draw[black!15] (11.25,1.0)--(11.25,9.05);
\node[font=\bfseries] at (5.5,8.95) {(a) 三个检测点与实际误差角域};
\node[font=\bfseries] at (17.3,8.95) {(b) 定位区域放大与覆盖检验};

% The full layout: true 2-degree sectors, no angular exaggeration.
\begin{scope}[shift={(5.65,7.25)},x=.009cm,y=.009cm]
  \clip (-585,-545) rectangle (540,125);
  \foreach \x in {-500,-250,0,250,500}
    \draw[black!7] (\x,-525)--(\x,110);
  \foreach \y in {-500,-250,0}
    \draw[black!7] (-560,\y)--(510,\y);
  \draw[->,black!35] (-560,0)--(510,0);
  \draw[->,black!35] (0,-525)--(0,110);

  \fill[qoneblue!15] (-500,0)--(400,0)--({-500+900*cos(2)},{900*sin(2)})--cycle;
  \fill[qoneorange!15] (-310,-450)
    --({-310+680*cos(atan(1.5)-2)},{-450+680*sin(atan(1.5)-2)})
    --({-310+680*cos(atan(1.5))},{-450+680*sin(atan(1.5))})--cycle;
  \fill[qonepurple!15] (310,-450)
    --({310+680*cos(180-atan(1.5))},{-450+680*sin(180-atan(1.5))})
    --({310+680*cos(182-atan(1.5))},{-450+680*sin(182-atan(1.5))})--cycle;

  \draw[qoneblue,thick] (-500,0)--(400,0);
  \draw[qoneblue,densely dashed] (-500,0)--({-500+900*cos(2)},{900*sin(2)});
  \draw[qoneorange,thick] (-310,-450)
    --({-310+680*cos(atan(1.5))},{-450+680*sin(atan(1.5))});
  \draw[qoneorange,densely dashed] (-310,-450)
    --({-310+680*cos(atan(1.5)-2)},{-450+680*sin(atan(1.5)-2)});
  \draw[qonepurple,thick] (310,-450)
    --({310+680*cos(180-atan(1.5))},{-450+680*sin(180-atan(1.5))});
  \draw[qonepurple,densely dashed] (310,-450)
    --({310+680*cos(182-atan(1.5))},{-450+680*sin(182-atan(1.5))});
  \filldraw[fill=qoneregion!65,draw=qoneregion,thick] (-10,0)--(10,0)--(0,15)--cycle;
  \draw[qoneink,densely dotted] (-28,-18) rectangle (28,35);

  \node[dot,fill=qoneblue] at (-500,0) {};
  \node[dot,fill=qoneorange] at (-310,-450) {};
  \node[dot,fill=qonepurple] at (310,-450) {};
\end{scope}
\node[align=left,anchor=west,qoneblue,fill=white,inner sep=2pt] at (0.6,7.9)
  {$S_1=(-500,0)$\\$\theta_1=1^\circ$};
\node[align=center,qoneorange] at (2.65,2.65)
  {$S_2=(-310,-450)$\\$\theta_2\approx55.309932^\circ$};
\node[align=center,qonepurple] at (8.3,2.65)
  {$S_3=(310,-450)$\\$\theta_3\approx124.690068^\circ$};
\draw[->,qoneink] (7.2,8.05)--(5.92,7.56);
\node[anchor=west,info] at (7.05,8.05) {公共区域 $P$};
\node[align=center,black!70,font=\footnotesize] at (5.5,1.65)
  {每个色带的张角均为 $2^\circ$，按实际比例绘制\\
   实线：形成三角形的边界\quad 虚线：不再裁剪三角形的边界};
\draw[->,black!40,thick] (10.4,7.45)--(11.8,7.45);
\node[font=\footnotesize,black!60,above] at (11.1,7.45) {放大};

% Zoom: the exact feasible triangle, diameter circle.
\begin{scope}[shift={(17.25,4.65)},scale=.21]
  \foreach \x in {-10,0,10} \draw[black!7] (\x,-12)--(\x,17);
  \foreach \y in {0,5,10,15} \draw[black!7] (-16,\y)--(16,\y);
  \draw[->,black!35] (-17,0)--(18,0) node[right] {$x$};
  \draw[->,black!35] (0,-12)--(0,18) node[above right] {$y$};
  \fill[qoneregion!12] (-10,0)--(10,0)--(0,15)--cycle;
  \begin{scope}
    \clip (-10,0)--(10,0)--(0,15)--cycle;
    \fill[qonefail!22,even odd rule]
      (-12,-12) rectangle (12,17) (0,0) circle[radius=10cm];
  \end{scope}
  \draw[qonefail,very thick] (0,0) circle[radius=10cm];
  \draw[qoneblue,very thick] (-10,0)--(10,0);
  \draw[qoneorange,very thick] (-10,0)--(0,15);
  \draw[qonepurple,very thick] (10,0)--(0,15);

  \node[dot] at (-10,0) {};
  \node[dot] at (10,0) {};
  \node[dot,fill=qonefail] at (0,15) {};
  \node[dot,fill=qonefail] at (0,0) {};

  \node[below left,info] at (-10,-.3) {$A=(-10,0)$};
  \node[below right,info] at (10,-.3) {$B=(10,0)$};
  \node[above left,info] at (-.5,15.2) {$U=(0,15)$};
  \node[below right,fill=white,inner sep=1pt,text=qonefail] at (.3,-.7) {$C_0$};
  \node[qoneregion] at (-4,4.6) {$P$};
  \node[text=qonefail,info] at (7,-5) {$r=10$};

  \draw[qonefail,densely dotted] (0,10)--(14.5,10);
  \draw[qonefail,densely dotted] (0,15)--(14.5,15);
  \draw[<->,qonefail,thick] (13.7,10)--(13.7,15)
    node[midway,right,fill=white,inner sep=2pt] {漏出 $5$ m};

  \draw[black!35] (-10,-9)--(-10,-12.2);
  \draw[black!35] (10,-9)--(10,-12.2);
  \draw[<->,qoneink] (-10,-11.6)--(10,-11.6)
    node[midway,below,fill=white,inner sep=2pt] {$D=|AB|=20$ m};
\end{scope}
\node[text=qonefail,font=\footnotesize,align=center] at (17.35,1.3)
  {红色实线：直径圆\quad $C_0=(0,0)$，$r=10$ m};
\end{tikzpicture}
}}

\usepackage[unicode,hidelinks]{hyperref}
\usepackage{bookmark}
\XeTeXgenerateactualtext=1
\hypersetup{pdftitle={无线电干扰源的快速自动定位与清除},pdfauthor={}}
\ctexset{section={format=\large\bfseries},subsection={format=\normalsize\bfseries},subsubsection={format=\normalsize\bfseries}}
\linespread{1.16}
\setlength{\parskip}{0.2em}
\setlength{\emergencystretch}{2em}
\setlist[enumerate]{label=\arabic*.,leftmargin=2em,itemsep=0.15em,topsep=0.4em,parsep=0pt}
\setlist[itemize]{leftmargin=2em,itemsep=0.15em,topsep=0.4em,parsep=0pt}
\pagestyle{fancy}\fancyhf{}\fancyfoot[C]{\thepage}
\renewcommand{\headrulewidth}{0pt}
\renewcommand{\footrulewidth}{0pt}
\widowpenalty=10000\clubpenalty=10000\displaywidowpenalty=10000
\raggedbottom
\begin{document}
\setcounter{page}{1}
\begin{center}{\LARGE\bfseries 无线电干扰源的快速自动定位与清除\par}\end{center}
\begin{abstract}
\noindent 摘要待四问模型与正式测试结果齐备后统一撰写。
\par\medskip\noindent\textbf{关键词：}交会定位；最小包围圆；主动测向；干扰源清除
\end{abstract}
\clearpage
\section{问题重述}
无线电干扰源分布在半径为 $1800$ m 的圆形区域内，其数量、位置和有效接收半径未知。机器狗通过带有 $\pm1^\circ$ 误差的示向度逐步定位，并在距目标不超过 $20$ m 时进行光学定位和清除。本文依次研究交会定位区域直径、第二检测点选择、全向源搜索清除以及包含未知朝向定向源时的策略。

\section{问题分析}
\subsection{问题一的分析}
问题一要求根据若干检测点的位置与示向度确定干扰源的可能区域，计算区域的直径，并判断以该直径为直径的圆能否覆盖整个区域。

由于测向存在 $\pm 1^\circ$ 的误差，一次测向只能把干扰源限制在一个张角 $2^\circ$ 的角域内，多次测向的定位区域即各角域的交集。角域可以写成两个线性不等式，于是定位区域是有限个\textbf{半平面的交}，非空有界时为凸多边形，其直径在顶点对上取得，枚举顶点即可求出。至于覆盖性，我们发现直径为 $D$ 的圆若要盖住区域，圆心只能取最远点对的中点，据此构造出一个符合题设测向条件的算例，其中该圆恰好盖不住区域的一个顶点，从而对第二小问给出否定回答。

\subsection{问题二的分析}
首测只能约束干扰源的方向与距离范围。第二点的选择需要同时解决两个问题：移动后仍能接收到信号，以及两次观测能够有效缩小定位区域。为此，先利用最小接收半径构造保证再接收的候选域，再结合交会角、可能源位置和移动代价选择第二点，最终用真实反馈计算更新后的区域。

\section{模型假设}
\begin{enumerate}
\item 测向误差有界，界为 $\pm 1^\circ$，不对误差的分布形式作附加假设；
\item 检测点位置精确已知，问题在平面内讨论；
\item 各测向数据相容，即真实干扰源同时位于所有误差角域内。
\item 问题二中的干扰源为全向源，其位置和有效接收半径在一次定位过程中保持不变，有效接收半径位于 $[1000,1500]$ m；
\item 同一地点的测向误差固定，重复测量不视为独立样本；机器狗可移动到目标圆域外的检测位置。
\end{enumerate}

\section{符号说明}
\begin{center}
\renewcommand{\arraystretch}{1.3}
\begin{tabularx}{0.98\linewidth}{>{\centering\arraybackslash}p{0.25\linewidth}>{\centering\arraybackslash}p{0.24\linewidth}>{\centering\arraybackslash}X}
\toprule
\textbf{符号} & \textbf{范围} & \textbf{说明} \\
\midrule
$n$ & $\mathbb{Z}^+$ & 测向次数 \\
$S_i=(x_i,y_i)$ & $\mathbb{R}^2$ & 第 $i$ 个检测点的坐标 \\
$\theta_i$ & $[0^\circ,360^\circ)$ & 第 $i$ 次测向的示向度 \\
$\varepsilon$ & $1^\circ$ & 测向误差界 \\
$P$ & --- & 定位区域 \\
$v_1,\dots,v_h$ & $\mathbb{R}^2$ & 定位区域的顶点 \\
$D$ & $[0,+\infty)$ & 定位区域的直径 \\
$C^*,\ R^*$ & --- & 最小包围圆的圆心与半径 \\
$\delta$ & $\pi\varepsilon/180$ & 用于三角函数计算的弧度误差界 \\
$\rho$ & $[1000,1500]$ m & 未知有效接收半径 \\
$\mathcal C_1,\ P_1$ & --- & 首测精确可行集及其外包多边形 \\
$K_0,\ Q$ & --- & 保证再接收域及有限候选检测点集 \\
$H$ & --- & 评分使用的有限假设源位置集 \\
$\lambda$ & $0.02$ m/s & 检测点评分中的移动时间权重 \\
\bottomrule
\end{tabularx}
\end{center}

\section{模型的建立与求解}
% Updated from current Markdown; old standalone TeX is preserved.
\subsection{问题一模型的建立与求解}\label{sec:q1}
\subsubsection{从测向数据到定位区域}\label{sec:q1-1}

设机器狗在 $n$ 个位置 $S_i=(x_i,y_i)$ 处对该干扰源测得示向度 $\theta_i$。$\theta_i$ 是自正东逆时针量到“检测点指向干扰源”方向的角度，但它并不是干扰源的真实方位：真实方位落在 $\theta_i$ 两侧各 $1^\circ$ 内。因此，干扰源一定位于以 $S_i$ 为顶点、方向在 $\theta_i\pm1^\circ$ 之间的一个 $2^\circ$ 窄楔形内。

现在要把“点在楔形内”写成数学不等式。记 $\varphi_i=\pi\theta_i/180$，$\delta=\pi\varepsilon/180=\pi/180$，楔形的两条边方向分别为

\[
u_i^-=\begin{pmatrix}\cos(\varphi_i-\delta)\\ \sin(\varphi_i-\delta)\end{pmatrix},\qquad
u_i^+=\begin{pmatrix}\cos(\varphi_i+\delta)\\ \sin(\varphi_i+\delta)\end{pmatrix},
\]

其中 $u_i^-$ 对应偏转角度的下限 $\varphi_i-\delta$，$u_i^+$ 对应上限 $\varphi_i+\delta$。判断一个点在两条边的哪一侧，可以用二维向量的叉积：对 $a=(a_x,a_y)$、$b=(b_x,b_y)$，定义 $a\times b=a_xb_y-a_yb_x$，它大于零说明 $b$ 在 $a$ 的逆时针一侧，小于零说明在顺时针一侧。于是，点 $X=(x,y)$ 落在楔形内的条件是

\begin{equation}\label{eq:q1-1}
u_i^-\times(X-S_i)\ge 0,\qquad u_i^+\times(X-S_i)\le 0.
\end{equation}

第一式要求 $X-S_i$ 不越过下边 $u_i^-$（在其逆时针一侧），第二式要求不越过上边 $u_i^+$（在其顺时针一侧），两式合起来就是把 $X$ 夹在两条边之间。

将式\eqref{eq:q1-1} 展开：

\begin{equation}\label{eq:q1-2}
\begin{aligned}
\sin(\varphi_i-\delta)(x-x_i)-\cos(\varphi_i-\delta)(y-y_i)&\le 0,\\
-\sin(\varphi_i+\delta)(x-x_i)+\cos(\varphi_i+\delta)(y-y_i)&\le 0.
\end{aligned}
\end{equation}

这两个式子都是 $x$、$y$ 的一次不等式，形如 $ax+by\le c$。满足 $ax+by\le c$ 的点集是直线 $ax+by=c$ 一侧的半个平面，称为半平面，它是一个凸集。式\eqref{eq:q1-2} 的两个不等式分别对应楔形的两条边，它们的交集就是该楔形。这样就把几何上“点在楔形内”的条件，转化成了代数上两个线性不等式同时成立。

$n$ 次测向后，干扰源必须同时落在每个楔形内，也就是要同时满足全部 $2n$ 个不等式。将它们统一写成 $a_jx+b_jy\le c_j$，得到定位区域

\begin{equation}\label{eq:q1-3}
P=\bigcap_{j=1}^{2n}\left\{(x,y):a_jx+b_jy\le c_j\right\}.
\end{equation}

$P$ 是 $2n$ 个半平面的交。半平面是凸集，凸集之交仍是凸集，故 $P$ 为凸集；在非空且有界时，$P$ 是\textbf{凸多边形}，退化情形为线段或点。题目图 2 中的红色四边形就是 $n=2$ 时两个楔形求交的结果。每增加一次有效测向，约束增加两条，定位区域随之收缩。

\subsubsection{求定位区域的直径}\label{sec:q1-2}

定位区域 $P$ 是凸多边形，现在要求它的直径---区域内距离最远的一对点之间的距离。

\textbf{首先确认区域有界。} 并非任意一组测向都能定出有限的区域。若各检测点近似共线、示向度方向接近平行，楔形之交沿某方向无限延伸，区域无界，直径为无穷大；若某次读数偏差较大，各楔形可能互不相容，区域为空。这两种情形都需要先排除。

具体做法是用线性规划：先检验约束组 $\{a_jx+b_jy\le c_j\}$ 是否可行，不可行说明数据矛盾，应剔除异常的测向数据；若可行，再分别求 $x$、$-x$、$y$、$-y$ 在约束下的最大值，四者均有限当且仅当 $P$ 有界。注意，无界区域也可能有若干有限顶点（例如楔形的尖端），若跳过这一步直接取顶点对距离，会把本应为无穷的直径误报成有限值。

\textbf{确认有界后，求出全部顶点。} 凸多边形的每个顶点都是某两条边界直线的交点。对直线 $a_jx+b_jy=c_j$ 与 $a_kx+b_ky=c_k$，记 $\Delta_{jk}=a_jb_k-a_kb_j$，当 $\Delta_{jk}\ne0$ 时交点为

\begin{equation}\label{eq:q1-4}
x_{jk}=\frac{c_jb_k-b_jc_k}{\Delta_{jk}},\qquad
y_{jk}=\frac{a_jc_k-c_ja_k}{\Delta_{jk}}.
\end{equation}

$\Delta_{jk}=0$ 时两直线平行，没有孤立交点，跳过。求出的交点未必在区域内，需代回全部 $2n$ 个约束检验，保留满足所有不等式者；去重后即得顶点集 $\{v_1,\dots,v_h\}$。

\textbf{有了顶点，直径就可以确定。} 按题目定义，$D=\max_{X,Y\in P}\|X-Y\|_2$。因为 $P$ 是凸多边形，我们证明这个最大值一定在顶点对上取得：

\begin{equation}\label{eq:q1-5}
D=\max_{1\le i<j\le h}\|v_i-v_j\|_2 .
\end{equation}

理由如下。凸多边形是其全部顶点的凸包，因此 $P$ 中任意一点都可写为顶点的凸组合：存在 $\lambda_i\ge0$、$\sum_i\lambda_i=1$，使 $X=\sum_i\lambda_iv_i$；同理 $Y=\sum_j\mu_jv_j$，$\mu_j\ge0$、$\sum_j\mu_j=1$。两式相减得

\[
X-Y=\sum_{i,j}\lambda_i\mu_j(v_i-v_j),
\qquad
\sum_{i,j}\lambda_i\mu_j=\left(\sum_i\lambda_i\right)\left(\sum_j\mu_j\right)=1 .
\]

由三角不等式，

\[
\|X-Y\|\le\sum_{i,j}\lambda_i\mu_j\|v_i-v_j\|\le\max_{i,j}\|v_i-v_j\|\cdot1=\max_{i,j}\|v_i-v_j\|,
\]

即区域内任意两点的距离都不超过最远顶点对的距离。而最远顶点对本身就是区域内的点，所以这个上界可以取到。式\eqref{eq:q1-5} 成立。

综上，求直径的完整流程为：

\begin{tcolorbox}[colback=black!2,colframe=black!22,boxrule=0.4pt,arc=0pt,left=9pt,right=9pt,top=7pt,bottom=7pt]
\small\setlength{\parskip}{2pt}
\noindent\hangindent=1.5em\hangafter=1 输入：检测点 $S_i$、示向度 $\theta_i$，$i=1,\dots,n$\par
\noindent\hangindent=1.5em\hangafter=1 1  按式\eqref{eq:q1-2} 将每次测向写成两个半平面约束；\par
\noindent\hangindent=1.5em\hangafter=1 2  线性规划检验可行性与有界性；\par
\noindent\hangindent=1.5em\hangafter=1 3  按式\eqref{eq:q1-4} 求每对边界直线的交点；\par
\noindent\hangindent=1.5em\hangafter=1 4  保留满足全部约束的交点并去重，得顶点集；\par
\noindent\hangindent=1.5em\hangafter=1 5  比较所有顶点对的距离，取最大者 $(A,B)$，$D=|AB|$。\par
\end{tcolorbox}

边界直线共 $2n$ 条，交点约 $\binom{2n}{2}=O(n^2)$ 个，每个交点检验约束需 $O(n)$，顶点数不超过 $2n$，比较顶点对为 $O(n^2)$，总计算量 $O(n^3)$。本题 $n$ 一般不超过十几，计算量完全可以实时完成。

\subsubsection{直径圆能否覆盖定位区域}\label{sec:q1-3}

机器狗的光学探测仪只能在距干扰源 $20$ m 以内工作。如果以定位区域直径 $D$ 为直径画一个圆，能盖住整个区域，那么只要 $D\le40$ m，机器狗走到圆心就能保证清除成功---这比计算最小包围圆要简单得多。所以问题是：这样的圆真的能盖住整个区域吗？

首先要确定圆心。设最远的一对顶点为 $A$、$B$，距离为 $D$。若半径 $D/2$、圆心 $C$ 的圆覆盖了 $P$，则 $A$、$B$ 都在圆内，即 $\|A-C\|\le D/2$、$\|C-B\|\le D/2$。代入三角不等式

\[
D=\|A-B\|\le\|A-C\|+\|C-B\|\le\frac D2+\frac D2=D .
\]

不等式的两端相等，故中间每一步都取等号：$\|A-C\|=\|C-B\|=D/2$，且 $C$ 在线段 $AB$ 上。这两个条件合起来，$C$ 只能是 $AB$ 的中点 $C_0=(A+B)/2$。圆是凸集，覆盖全部顶点就覆盖了整个凸多边形，所以判据为

\begin{equation}\label{eq:q1-6}
\max_{1\le i\le h}\|v_i-C_0\|\le\frac D2 .
\end{equation}

式\eqref{eq:q1-6} 是否一定满足？下面说明它可以不成立，并给出一个满足题设条件的具体反例。

\textbf{构造思路。} 要让直径圆盖不住区域，需要找到这样一个顶点：它已经跑到了直径圆外面，但它到其他顶点的距离还没有超过原来的直径---直径没变，圆却已经不够大了。

先用归一化的坐标把这个想法量化。取 $A=(-1,0)$、$B=(1,0)$，$|AB|=2$，直径圆以原点为圆心、半径为 $1$。在上方放一个点 $U=(0,h)$，希望它同时满足：

\begin{itemize}
\item 跑到圆外：$h>1$；
\item 不改变直径：$|UA|=|UB|=\sqrt{1+h^2}\le2$，即 $h\le\sqrt{3}$。
\end{itemize}

两个条件合起来是 $1<h\le\sqrt{3}$，这个区间非空。直观地理解：上顶点刚越过圆的最高点时，它离左右两端还不够远，不会产生更长的边。

取 $h=1.5$，则 $|UA|=\sqrt{3.25}\approx1.803<2$，直径仍为 $2$；而上顶点距圆心 $1.5$，超出半径 $50\%$。反例的形状就有了：底边为 $2$、高为 $1.5$ 的等腰三角形。

\textbf{用测向数据实现这个形状。} 将坐标放大到米，目标三角形的三个顶点为

\[
A=(-10,\,0),\quad B=(10,\,0),\quad U=(0,\,15).
\]

三条边分别是 $y=0$（底边）、$y=1.5x+15$（左腰）、$y=-1.5x+15$（右腰）。三角形有 $3$ 条边，而每个检测点的楔形提供 $2$ 条边界线，因此用 $3$ 个检测点（$6$ 条边界线）即可实现：让其中 $3$ 条边界线分别与三角形的三条边重合，另外 $3$ 条不切入三角形内部。具体取法见表\ref{tab:q1-counterexample}。

\Needspace{19\baselineskip}
\begin{center}
\captionof{table}{反例的检测数据（误差界 $\pm1^\circ$）}\label{tab:q1-counterexample}
\end{center}

\begin{center}
\renewcommand{\arraystretch}{1.3}
\begin{tabularx}{0.98\linewidth}{>{\centering\arraybackslash}p{0.22\linewidth}>{\centering\arraybackslash}p{0.20\linewidth}>{\raggedright\arraybackslash}X}
\toprule
\textbf{检测点 $S_i$（m）} & \textbf{示向度 $\theta_i$} & \textbf{起决定作用的边界} \\
\midrule
$(-500,\,0)$ & $1^\circ$ & 下边界方向为 $0^\circ$，给出 $y\ge0$（底边） \\
$(-310,\,-450)$ & $\approx55.31^\circ$ & 上边界方向为 $\arctan(1.5)$，给出 $y\le1.5x+15$（左腰） \\
$(310,\,-450)$ & $\approx124.69^\circ$ & 下边界方向为 $180^\circ-\arctan(1.5)$，给出 $y\le-1.5x+15$（右腰） \\
\bottomrule
\end{tabularx}
\end{center}

其中第二个检测点的示向度取 $\theta_2=\arctan(1.5)-1^\circ\approx55.31^\circ$，第三个取 $\theta_3=181^\circ-\arctan(1.5)\approx124.69^\circ$。将三角形的三个顶点分别代入每个楔形的另一条边界线，可以验证它们均满足约束（不被裁掉），因此 $6$ 个半平面的交恰好就是三角形 $ABU$。

现在检验覆盖性。边长为 $|AB|=20$，$|AU|=|BU|=\sqrt{325}\approx18.03$，直径 $D=20$ m 由 $A$、$B$ 取得。直径圆圆心 $C_0=(0,\,0)$，半径 $D/2=10$ m。而上顶点到圆心的距离为

\[
\|U-C_0\|=15\ \text{m}>10\ \text{m}=\frac D2,
\]

超出半径 $5$ m，式\eqref{eq:q1-6} 不成立。

% Added case illustration; the six constraints use the exact stated angles.
\begin{center}
\begin{minipage}{\linewidth}
\centering
\QOneCaseFigure
\captionof{figure}{三次测向的角域交集与直径圆覆盖检验。左图按实际比例绘制三个误差角域；右图放大定位三角形，红色圆的半径为 $10$ m，上顶点超出 $5$ m。}
\label{fig:q1-counterexample}
\end{minipage}
\end{center}


三个检测点到三角形中心的距离均约 $500$ m，在有效接收半径内；三角形整体位于半径 $1800$ m 的目标区域内，是完全合理的探测场景。于是问题一的第二问答案为否：

\textbf{以定位区域直径为直径的圆，不一定能覆盖该定位区域。}

% Second-question draft. Geometry and parameters follow repository main.
\subsection{问题二模型的建立与求解}\label{sec:q2}

\subsubsection{首测信息与候选源区域}

问题一将多次示向度转化为定位区域，并给出了最小包围圆半径 $R^*$ 对位置不确定性的刻画。问题二进一步要求：在只有首测信息时，主动选择第二检测点，使再次测向后的位置不确定性尽量减小。由于首测只提供方向区间而不提供距离，不能把示向射线上的某一点当作真实源位置，也不能直接按未知源坐标构造直角交会。

沿用第一问记号，首个检测点为 $S_1$，示向度为 $\theta_1$，对应误差角域为 $W_1$。以 $\overline B(c,r)$ 表示闭圆盘，目标区域为 $\Omega=\overline B((0,0),1800)$。普通示向响应说明源与首站距离大于 $5$ m，并且不超过有效接收半径的上限 $1500$ m，因此首测精确可行集为
\begin{equation}\label{eq:q2-source}
\mathcal C_1=\bigl(\Omega\cap W_1\cap\overline B(S_1,1500)\bigr)\setminus\overline B(S_1,5).
\end{equation}
为使用第一问的凸多边形运算，计算时省略 $5$ m 内圆排除，并用外切正多边形包络圆边界，得到 $P_1\supseteq\mathcal C_1$。这里 $P_1$ 表示源可能位于何处，下面的 $K_0$ 则表示第二检测点可以选在何处，二者的几何含义不同。

\subsubsection{保证再次接收的第二点候选区域}

全向源的有效接收半径 $\rho\in[1000,1500]$ m 未知。为避免移动后完全失去信号，先要求第二点 $q$ 到每个候选源位置都不超过最小接收半径，定义
\begin{equation}\label{eq:q2-safe}
\begin{aligned}
K_0(P_1)&=\left\{q:\max_{X\in P_1}\|q-X\|_2\le1000\right\}\\
&=\bigcap_{v\in\operatorname{Vert}(P_1)}\overline B(v,1000).
\end{aligned}
\end{equation}
第二个等号不需要遍历整个区域。事实上，对任意 $X=\sum_i\alpha_i v_i\in P_1$，其中 $\alpha_i\ge0$ 且 $\sum_i\alpha_i=1$，有
\[
\|q-X\|_2\le\sum_i\alpha_i\|q-v_i\|_2\le\max_i\|q-v_i\|_2.
\]
因此只要逐顶点距离检查通过，就能覆盖所有候选位置。又因真实源位于 $P_1$ 内且 $\rho\ge1000$ m，任意 $q\in K_0$ 均保证再次接收到信号：距离大于 $5$ m 时返回示向度，不超过 $5$ m 时返回近场响应，可直接进行光学定位。

第一问的最小包围圆还给出候选域的非空判据：
\begin{equation}\label{eq:q2-nonempty}
K_0(P_1)\ne\varnothing\quad\Longleftrightarrow\quad R^*(P_1)\le1000.
\end{equation}
此时 $C^*(P_1)$ 就是一个可行第二点。为直观说明首测后存在可选空间，令 $u(\varphi_1)=(\cos\varphi_1,\sin\varphi_1)$，完整首测扇形可被下述圆包络：
\[
c_1=S_1+\frac{1500}{2\cos\delta}u(\varphi_1),\qquad R_1=\frac{1500}{2\cos\delta}.
\]
沿任一射线，到 $c_1$ 的距离平方关于径向距离为凸函数，检查扇形顶点和外弧端点即可验证包络。当 $\delta=\pi/180$ 时 $R_1\approx750.114$ m，故 $\overline B(c_1,1000-R_1)$ 给出一个明确的保证再接收区域；加入目标圆先验只会减小源候选集。实际离散计算仍以式\eqref{eq:q2-safe}逐顶点核验为准。

$K_0$ 是\textbf{保守的保证再接收域}。首测有信号还说明 $\rho\ge\|X-S_1\|$，充分利用这一联合信息可能允许更多第二点；本文选用易于验证的式\eqref{eq:q2-safe}，不将它称为最大可行域。实现中将各半径 $1000$ m 圆盘替换为内接多边形以生成候选；内近似为空不能单独据此判定真实圆弧交集为空，仍须使用式\eqref{eq:q2-nonempty}检验。

\subsubsection{交会几何与候选点评分}

再次接收只是可行性要求，还需避免两次视线接近平行。对假设源位置 $X$，记两站到源的距离为 $r_1=\|X-S_1\|$、$r_2=\|X-q\|$，两条视线的锐交会角为 $\psi$。对测向方程作局部线性化，横向位置误差的尺度约为 $r_i\tan\delta$，两条约束法向量所组成矩阵的行列式绝对值为 $|\sin\psi|$，由此得到误差代理尺度
\begin{equation}\label{eq:q2-error}
E(q,X)=\frac{(r_1+r_2)\tan\delta}{|\sin\psi(q,X)|}.
\end{equation}
该式解释了交会角较大通常有利、近平行方向误差放大的原因，但它是一阶近似，并不是原非线性定位区域半径的严格上界。未知 $X$ 时，应对多个相容位置评价，而非只对首测射线上的一个估计点设计交会。

根据首测可得的 $P_1$ 和 $K_0$ 构造有限候选检测点集 $Q$，同时从 $P_1$ 的顶点及最小圆心构造有限假设集 $H$。采用定位误差代理与移动时间的组合评分
\begin{equation}\label{eq:q2-score}
J(q)=\max_{X\in H}\frac{(\|X-S_1\|+\|X-q\|)\tan\widetilde\delta}
{\max\{|\sin\psi(q,X)|,0.02\}}+\lambda\frac{\|q-S_1\|}{5},
\end{equation}
其中速度为 $5$ m/s，$\lambda=0.02$ m/s，使两项均具有长度量纲。$\widetilde\delta=1.005\pi/180$ 是实现中使用的保守半角；相对于题面 $1^\circ$，额外的 $0.005^\circ$ 用于容纳接口两位小数示向的舍入，构造 $P_1$ 时使用同一半角。分母下限 $0.02$ 用于控制近平行情况的数值发散。这些是评分与实现参数，不是新的物理假设。

如果假设源距第二点不超过 $5$ m，评分使用近场尺度 $5$ m；外包角域尖点引入的零首测距离假设不参与角度计算。同地点误差固定，故在生成 $Q$ 时剔除距已有检测站过近的候选，默认最小间距取 $25$ m。

\subsubsection{选点算法及第二次反馈处理}

具体步骤如下，所用信息均来自题设先验和已经取得的示向度。
\begin{enumerate}
\item 计算首测外包多边形 $P_1$ 及其最小包围圆 $(C^*,R^*)$，按式\eqref{eq:q2-nonempty}检验保证再接收域是否非空。
\item 构造 $K_0$ 的安全内多边形，候选中加入 $C^*$、至多 $16$ 个内多边形顶点，以及 $C^*\pm h(-\sin\varphi_1,\cos\varphi_1)$，其中 $h\in\{50,150,300,500\}$ m；去重并剔除距首站不足 $25$ m 的点。
\item 对每个候选重新检查到 $P_1$ 全部顶点的距离，保留满足式\eqref{eq:q2-safe}者；用至多 $12$ 个抽样源域顶点及最小圆心组成 $H$。
\item 计算 $J(q)$，选择有限集合 $Q$ 中评分最小的检测点；并列时依次选择移动距离较短、坐标字典序较小的点。
\item 在选定点测量。若返回示向度，则将第二次误差角域、接收距离上界与原区域求交，重新计算 $D$ 和 $R^*$；若返回近场响应，则在当前点进行光学定位。后验区域是否足够小以实际反馈为准，不由预测评分代替。
\end{enumerate}
该算法保证保留候选的再次接收，并在\textbf{有限候选集内}择优。顶点抽样、误差线性化和分母截断使它不构成连续检测域上的全局最优算法。若程序因数据不一致或离散筛选未获得合格候选，应报告这一状态，不能把任意位置标记为保证接收点。

\subsubsection{选点算例与定位效果}

取 $S_1=(0,0)$、首测 $\theta_1=0^\circ$。选点算法只使用首测信息；为比较第二次测量后的区域，另取真实源 $G=(1000,0)$ 构造一个受控算例。$G$ 仅用于生成观测及事后检验，没有传入选点函数。第二次示向采用无附加误差的真实方位并舍入至两位小数，区域计算采用 $1.005^\circ$ 保守半角。

\begin{table}[htbp]
\centering
\caption{首测条件相同的第二检测点对照}\label{tab:q2-example}
\small
\begin{tabular}{lccc}
\toprule
方案 & 第二检测点 / m & 第二示向 & 后验 $R^*$ / m \\
\midrule
仅有首测 & --- & --- & 750.2308 \\
共线前进 & $(500,0)$ & $0.00^\circ$ & 500.1539 \\
主动选点 & $(857.728610,514.102744)$ & $285.47^\circ$ & 22.8123 \\
\bottomrule
\end{tabular}
\end{table}

本例得到 $25$ 个合格候选，所选点通过全部顶点的 $1000$ m 距离核验。表\ref{tab:q2-example}表明，共线前进仍保留较大的纵向不确定性，而横向选点显著改善交会几何。两方案分别移动 $500$ m 和 $1000$ m，故此表用于说明几何缩域效果，不能据此宣称等成本最优或完整清除任务的时间优势。

\begin{figure}[htbp]
\centering
\begin{tikzpicture}[x=0.006cm,y=0.006cm,font=\small]
\draw[->,gray] (-100,0)--(1650,0) node[right] {$x$/m};
\draw[->,gray] (0,-740)--(0,800) node[above] {$y$/m};
\filldraw[fill=green!12,draw=green!45!black] (998.795456,-49.067674) -- (998.795456,49.067674) -- (989.176510,146.730474) -- (970.031253,242.980180) -- (941.544065,336.889853) -- (903.989293,427.555093) -- (857.728610,514.102744) -- (803.207531,595.699304) -- (760.969704,647.166247) -- (759.048875,645.245418) -- (696.792469,569.385767) -- (642.271390,487.789207) -- (596.010707,401.241556) -- (558.455935,310.576316) -- (529.968747,216.666643) -- (510.823490,120.416937) -- (501.204544,22.754137) -- (501.204544,-22.754137) -- (510.823490,-120.416937) -- (529.968747,-216.666643) -- (558.455935,-310.576316) -- (596.010707,-401.241556) -- (642.271390,-487.789207) -- (696.792469,-569.385767) -- (759.048875,-645.245418) -- (760.969704,-647.166247) -- (803.207531,-595.699304) -- (857.728610,-514.102744) -- (903.989293,-427.555093) -- (941.544065,-336.889853) -- (970.031253,-242.980180) -- (989.176510,-146.730474) -- cycle;
\filldraw[fill=blue!22,draw=blue!60!black,line width=.7pt] (1500.000000,26.313537) -- (0.000000,0.000000) -- (1500.000000,-26.313537) -- cycle;
\draw[densely dashed,->,red!70!black] (0,0)--(857.728610,514.102744);
\fill (0,0) circle[radius=1.7pt] node[below left] {$S_1$};
\fill[red!70!black] (857.728610,514.102744) circle[radius=1.8pt] node[above right] {$S_2$};
\fill (500,0) circle[radius=1.5pt] node[below] {共线对照点};
\fill[orange!80!black] (1000,0) circle[radius=1.6pt] node[above right] {$G$（仅评价）};
\node[green!30!black] at (750,-340) {$K_0$ 的内近似};
\node[blue!70!black] at (1370,170) {$P_1$};
\draw[blue!60!black,->] (1370,135)--(1370,23);
\foreach \x in {500,1000,1500}{\draw[gray] (\x,0)--(\x,-18);}
\end{tikzpicture}
\caption{首测源候选区域与第二检测点的安全候选域。坐标等比例，窄扇形按实际半角绘制；绿色区域为保证再接收域的内近似。}\label{fig:q2-domain}
\end{figure}

还应注意，主动选点后的 $R^*=22.8123$ m 仍大于 $20$ m，因此本例尚不能保证一次光学清除，需要继续观测或使用后续问题的有限覆盖清除策略。该算例也不意味着任意源位置和误差组合都能在两次测向后达到相同精度。问题二得到的是可执行的检测点候选域和选取方法，其有效性最终由实际第二次观测形成的定位区域检验。

% 问题三、四由后续章节接入；保留已有编号和标签。
% 参考文献、AI 工具使用声明、附录及支撑材料清单待全稿整合。
\end{document}
